% !PRINTED 230
\weberpage{230}

%% [running head, not body text -- commented out 2026-08-07] Page furniture (running
%% head, folio, signature) is excluded from this edition by standing policy.
%% Sechster Abschnitt. \hfill \S.\,74.

\begin{center}
\S.\,74.

\gesperrt{Transformation der Gleichung fünften Grades.}
\end{center}

\indent Diese Entwickelungen sollen nun angewandt werden, um
die Transformation der Gleichung fünften Grades, wie wir sie
im \S.\,54 skizzirt haben, durchzuführen.

\indent Es lässt sich auf unendlich viele Arten ein Werthsystem
$t_0, t_1, t_2, t_3$ so bestimmen, dass die Bezoutiante $B$ verschwindet;
im Allgemeinen ist dazu die Auflösung einer quadratischen
Gleichung erforderlich. Wir können z.\,B.
\[
t_2 = 0, \quad t_3 = 0
\]
setzen und das Verhältniss $t_0 : t_1$ aus der quadratischen Gleichung
\[
B_{0,0}\,t_0^2 + 2\,B_{0,1}\,t_0\,t_1 + B_{1,1}\,t_1^2 = 0\,.
\]
bestimmen. Diese Werthe von $t$ enthalten dann die Quadratwurzel
\[
\sqrt{B_{0,1}^2 - B_{0,0}\,B_{1,1}}\,.
\]
Statt dessen kann man aber auch andere Bestimmungen treffen
und bekommt andere und andere Quadratwurzeln. In der Auswahl
dieser Quadratwurzel liegt etwas Willkürliches und Unbestimmtes.

\indent Wir greifen unter diesen verschiedenen Bestimmungsweisen
eine heraus und transformiren damit die gegebene Gleichung
fünften Grades in eine \gesperrt{Hauptgleichung}, d.\,h. in eine solche,
in der die dritte und vierte Potenz der Unbekannten nicht vorkommt, oder wir nehmen an, nach F. Klein's Vorgang, die zu
transformirende Gleichung sei \gesperrt{von Haus aus eine Hauptgleichung}
\[
% The source's principal-form specialization intentionally has no x^4 or x^3 terms.
\text{(1)} \qquad f(x) = a_0\,x^5 + a_3\,x^2 + a_4\,x + a_5 = 0\,.
\]
Wir wollen aber noch ausdrücklich hervorheben, dass nach
\S.\,70, (8) durch diese vorläufige Tschirnhausen-Transformation
die Discriminante der gegebenen Gleichung nur um einen Factor
$\Theta^2$ geändert wird, worin $\Theta$ \gesperrt{rational} von den angewandten $t$,
also auch von der zu ihrer Bestimmung benutzten Quadratwurzel
abhängt.

\indent Unter der Voraussetzung (1) werden nun die Co\"efficienten
der Bezoutiante nach den Formeln des vorigen Paragraphen
folgende:
%CONT>>
